%=============================================================================
% EXPERIMENTAL SETUP — Dataset, Protocol, Metrics, Baselines
%
% Rewritten after the Phase-2 review:
%   - peak timing error redefined (the previous metric never looked along time)
%   - the dense fixed-graph control added to the baseline list
%   - reporting now states the aggregation explicitly and reports dispersion
% See docs/PHASE2_REVIEW.md F4, F6, F7, F9.
%=============================================================================
\section{Experimental Setup}
\label{sec:setup}

\subsection{Dataset}
We evaluate on Sri Lanka's national dengue surveillance record, $N=25$
administrative districts over $T=459$ consecutive weeks (2013--2022). Weekly case
counts come from the Ministry of Health Epidemiology Unit; ten environmental
covariates (precipitation, air and land-surface temperature, humidity, NDVI, soil
moisture, canopy cover) are drawn from NASA EarthData. Covariates are pre-shifted
by their empirically optimal lags --- precipitation $\sim$12 weeks, NDVI
$\sim$17 weeks --- reflecting vector breeding and infection-to-diagnosis
intervals. The target is strongly right-skewed: median 13 weekly cases, maximum
2{,}631, and $9.7\%$ of district-weeks zero.

\subsection{Evaluation Protocol}
Models are evaluated under leakage-free rolling-origin (expanding window)
cross-validation across three chronological origins
($a \in \{0.55, 0.70, 0.85\}$) and three seeds, giving nine runs per
configuration. Each origin uses the following $15\%$ of the series as test data
and the final 30 weeks before the origin as validation. Standardisation
statistics are computed strictly from training weeks. Each run makes one
optimiser step per training window.

A rolling-origin protocol is necessary rather than stylistic. Under a single
$70/10/20$ chronological split we measured validation and test RMSE to be
anti-correlated ($r = -0.74$): the two periods fall in different epidemic
regimes, so selecting on one actively mis-selects for the other.

Results are reported as mean $\pm$ standard deviation over the nine runs. RMSE is
pooled across horizons within a fold, then averaged across folds; we state this
because pooling and averaging per-horizon RMSE differ by roughly $0.8$ here,
larger than several of the effects below. Comparisons against persistence use a
paired sign test over matched $(\text{fold}, \text{seed})$ runs rather than a
comparison of means.

\subsection{Metrics}
Point accuracy is measured by RMSE, MAE and SMAPE on the original case scale;
SMAPE is the primary percentage metric because frequent zero-case weeks make
plain MAPE unusable.

We additionally report \emph{peak week error}: per district, the absolute
difference in weeks between the observed and forecast week of maximum incidence,
averaged over districts whose observed maximum reaches ten cases. Below that
threshold there is no outbreak to time and the argmax is noise. This is what an
early-warning system is judged on, and it is not recoverable from RMSE. As a
check, a persistence forecast at horizon $h$ is shifted $h$ weeks late by
construction and must score approximately $h$; it does
(Table~\ref{tab:horizon}).

\subsection{Baselines}
\textbf{(1) Persistence floor:} $\hat{y}_{t+h} = y_t$, deterministic, one run per
fold.
\textbf{(2) Dense GCN, fixed graph:} the matched control of
Section~\ref{sec:framework}, sharing the propagation path so any difference is
attributable to the learned adjacency alone.
We do not compare numerically against Weng et al.~\cite{weng2024graph}: their
static $70/30$ split and different target normalisation make the numbers
incommensurable.
